On the prespecified 36-task subset, the fixed interpreted portfolio produces 36 valid plans and meets 27 deadlines, at a measured mean preparation/check interval of 20.13\,ms. The native deadline model meets 31 deadlines. The solver reports four deadline models infeasible and reaches one solve limit without a qualifying result. Its mean native-and-check interval is 278.11\,ms; the practical workflow that also retains the known portfolio fallback averages 298.24\,ms and preserves 36 valid results.

The makespan model likewise qualifies 31/36 and returns 36 valid plans. It reports 29 optimal objective statuses and seven time-limited incumbents; the raw gap and dual-bound fields are preserved. Its mean end-to-end interval, including the independently checked portfolio upper bound, is 1,021.78\,ms. Shared native process startup adds 405.02\,ms to the campaign. These are measured physical intervals on one host, not conversions from the primary work tariff. A deadline mode's optimal status concerns its zero-objective feasibility model, not optimal makespan; returned feasible schedules are independently checked, while infeasibility and optimality statuses are solver reports without separately verified proof certificates.

The additional four qualified deadlines come from a strong established specialist. They are not improvements discovered by the flywheel. The reference illustrates a practical quality/latency choice that a future hybrid may exploit under a common, prospectively measured budget. The present experiment does not establish such a hybrid's overall superiority or price its extra computation as revenue.
